\chapter{\IfLanguageName{dutch}{Testing en evaluatie van het systeem}{Testing and Evaluation of the System}}%
\label{ch:testing-en-evaluatie-van-het-systeem}
In dit hoofdstuk wordt het ontwikkelde systeem als geheel getest en geëvalueerd. Hierbij wordt nagegaan of de verschillende onderdelen van de applicatie correct samenwerken en of OBD-II-voertuigdata succesvol kan worden uitgelezen, verwerkt en weergegeven in het realtime monitoringdashboard.
De applicatie wordt getest met echte voertuigdata en, waar relevant, met de ontwikkelde simulatieomgeving. De resultaten worden vervolgens gebruikt om te beoordelen in welke mate de ontwikkelde oplossing voldoet aan de vooropgestelde doelstellingen van het onderzoek.

\section{Functionele testen}
Tijdens de functionele testen wordt nagegaan of de volledige datastroom correct werkt. Hierbij wordt gecontroleerd of een verbinding met het voertuig kan worden opgebouwd, voertuigdata kan worden uitgelezen en de ontvangen gegevens door de backend kunnen worden verwerkt.
Vervolgens wordt nagegaan of de data correct naar de frontend wordt doorgestuurd en of de ontvangen waarden correct in het dashboard worden weergegeven. Ook wordt gecontroleerd of nieuwe voertuiggegevens realtime zichtbaar worden wanneer de waarden veranderen.

\section{Testen met echte voertuigdata}
De applicatie wordt getest met de voertuigen die eerder in het onderzoek werden gebruikt. Hierbij wordt nagegaan of de geselecteerde OBD-II-parameters correct kunnen worden uitgelezen en weergegeven in het dashboard.
De tests hebben voornamelijk als doel om te controleren of de volledige keten, van het voertuig en de OBD-II-adapter tot de backend en frontend, correct functioneert met echte voertuigdata.
\section{Testen met gesimuleerde data}
De simulatieomgeving wordt gebruikt om de applicatie te testen wanneer geen fysiek voertuig beschikbaar is. Hierbij wordt nagegaan of de gegenereerde voertuigdata op dezelfde manier door de backend en frontend kan worden verwerkt als echte OBD-II-data.
De simulator wordt voornamelijk geëvalueerd als hulpmiddel voor ontwikkeling en functionele testing en wordt niet gebruikt voor de performantievergelijking tussen REST API en WebSockets.

\section{Evaluatie en beperkingen}
Op basis van de uitgevoerde testen wordt beoordeeld in welke mate de ontwikkelde applicatie geschikt is voor realtime voertuigmonitoring. Hierbij wordt gekeken naar de correcte werking van de volledige datastroom, de bruikbaarheid van het dashboard en de inzetbaarheid van de simulator tijdens de ontwikkeling.
Het onderzoek kent enkele beperkingen. Zo werd slechts een beperkt aantal voertuigen onderzocht en kunnen de beschikbare OBD-II-parameters verschillen tussen voertuigen. Daarnaast vormt de simulator een vereenvoudigde weergave van echte voertuigdata en kan deze niet alle eigenschappen van een werkelijk voertuig nabootsen.
De resultaten uit dit hoofdstuk worden samen met de vergelijking tussen REST API en WebSockets gebruikt om de uiteindelijke werking van het ontwikkelde systeem te beoordelen.
